\section{Produktquadrate mit höchstens zwei Werten}
Die Potenzmenge des Produktquadrats darf im Allgemeinen nicht mit dem
Produktquadrat der Potenzhalbgruppe gleichgesetzt werden. Bei einer
Zweiermenge reichen jedoch die Einermengen und der ganze Träger aus.

\Needspace{14\baselineskip}
\FormulaThmDeltaK[Einermengen und Gesamtmenge im Potenzquadrat]{%
\begin{aligned}[t]
&\SemigroupAlg{A}{\star}\dsep A\neq\varnothing\vdash{}\\[-2pt]
&A^2\neq\varnothing\land\PowNE{A}^{\,2}\subseteq\PowNE{A^2}\\[-2pt]
&\qquad{}\land A^2\in\PowNE{A}^{\,2}
 \land\forall a\in A^2\,(\{a\}\in\PowNE{A}^{\,2})
\end{aligned}
}{SemigroupSquarePowerFamily}{%
  \DeltaRow{Mengen}{A\dsep a}
  \DeltaRow{Funktionen}{\star\dsep\star_{\mathcal P}}
}
\begin{tabproofwide}
  \proofstepwidestar[1]{\SemigroupAlg{A}{\star}}{\rA}
  \proofstepwidestar[2]{A\neq\varnothing}{\rA}
  \proofstepwidestar[2]{A\in\PowNE{A}}{\FormulaRefAuto{NonemptyPowersetMembership}[thm]{\rAI{\FormulaRefAuto{A\subseteq A}[thm],2}}}
  \proofstepwidestar[1,2]{A^2=A\star_{\mathcal P}A}{\FormulaRefAuto{SemigroupSquareDef}[def]{1,2}}
  \proofstepwidestar[1,2]{A\star_{\mathcal P}A\neq\varnothing}{\FormulaRefAuto{PowerSemigroupProductNonempty}[thm]{1,3,3}}
  \proofstepwidestar[1,2]{A^2\neq\varnothing}{\rIE{\FormulaRefAuto{a=b\vdash b=a}[thm]{4},5}}
  \proofstepwidestar[1,2]{\PowNE{A}^{\,2}\subseteq\PowNE{A^2}}{\FormulaRefAuto{PowerSemigroupSquareSubset}[thm]{1,2}}
  \proofstepwidestar[1]{\SemigroupAlg{\PowNE{A}}{\star_{\mathcal P}}}{\FormulaRefAuto{NonemptyPowerSemigroup}[thm]{1}}
  \proofstepwidestar[2]{\PowNE{A}\neq\varnothing}{\FormulaRefAuto{NonemptyPowerCarrierNonempty}[thm]{2}}
  \proofstepwidestar[1,2]{A\star_{\mathcal P}A\in\PowNE{A}^{\,2}}{\FormulaRefAuto{SemigroupSquareProductMember}[thm]{8,9,\rAI{3,3}}}
  \proofstepwidestar[1,2]{A^2\in\PowNE{A}^{\,2}}{\rIE{\FormulaRefAuto{a=b\vdash b=a}[thm]{4},10}}
  \proofstepwidestar[12]{a\in A^2}{\rA}
  \proofstepwidestar[1,2]{A^2\subseteq A}{\FormulaRefAuto{SemigroupSquareSubset}[thm]{1,2}}
  \proofstepwidestar[1,2,12]{a\in A}{\FormulaRefAuto{A\subseteq B\dsep x\in A\vdash x\in B}[thm]{13,12}}
  \proofstepwidestar[1,2,12]{\{a\}\in\PowNE{A}^{\,2}}{\rRE{\rLREb{\FormulaRefAuto{PowerSemigroupSquareSingletonCriterion}[thm]{1,2,14}},12}}
  \proofstepwidestar[1,2]{\forall a\in A^2\,(\{a\}\in\PowNE{A}^{\,2})}{\rUI{\rRI{12,15}}}
  \proofstepwidestar[1,2]{\begin{gathered}A^2\neq\varnothing\land\PowNE{A}^{\,2}\subseteq\PowNE{A^2}\\{}\land A^2\in\PowNE{A}^{\,2}\\
 {}\land\forall a\in A^2\,(\{a\}\in\PowNE{A}^{\,2})\end{gathered}}{\rAI{6,\rAI{7,\rAI{11,16}}}}
\end{tabproofwide}
\Needspace{14\baselineskip}
\FormulaThmDeltaK[Bei einem zweielementigen Quadrat gibt es keine fehlenden Teilmengen]{%
\begin{aligned}[t]
&\SemigroupAlg{A}{\star}\dsep A\neq\varnothing\dsep A^2=\{a,b\}\\[-2pt]
&\vdash\PowNE{A}^{\,2}=\PowNE{A^2}
\end{aligned}
}{PairProductSquarePowerEquality}{%
  \DeltaRow{Mengen}{A\dsep a\dsep b\dsep X}
  \DeltaRow{Funktionen}{\star\dsep\star_{\mathcal P}}
}
Der Beweis gilt auch für \(a=b\). In der Tabelle steht
\(D=A^2\) und \(K=\PowNE{A}^{\,2}\); dies sind ausschließlich Abkürzungen.

\begin{tabproofwide}
  \proofstepwidestar[1]{\SemigroupAlg{A}{\star}}{\rA}
  \proofstepwidestar[2]{A\neq\varnothing}{\rA}
  \proofstepwidestar[3]{D=\{a,b\}}{\rA}
  \proofstepwidestar[1,2]{\begin{gathered}D\neq\varnothing\land K\subseteq\PowNE{D}\\{}\land D\in K\land\forall d\in D\,(\{d\}\in K)\end{gathered}}{\FormulaRefAuto{SemigroupSquarePowerFamily}[thm]{1,2}}
  \proofstepwidestar[3]{a\in D}{\rIE{\FormulaRefAuto{a=b\vdash b=a}[thm]{3},\FormulaRefAuto{a\in\{a,b\}}[thm]}}
  \proofstepwidestar[3]{b\in D}{\rIE{\FormulaRefAuto{a=b\vdash b=a}[thm]{3},\FormulaRefAuto{b\in\{a,b\}}[thm]}}
  \proofstepwidestar[1,2,3]{\{a\}\in K}{\rRE{5,\rUE{\rAEn{4}}}}
  \proofstepwidestar[1,2,3]{\{b\}\in K}{\rRE{6,\rUE{\rAEn{4}}}}
  \proofstepwidestar[1,2,3]{\{a,b\}\in K}{\rIE{3,\rAEn{4}}}
  \proofstepwidestar[10]{X\in\PowNE{D}}{\rA}
  \proofstepwidestar[3,10]{X\in\PowNE{\{a,b\}}}{\rIE{3,10}}
  \proofstepwidestar[3,10]{X\in\bigl\{\{a\},\{b\},\{a,b\}\bigr\}}{\rIE{\FormulaRefAuto{NonemptyPowersetPair}[thm],11}}
  \proofstepwidestar[3,10]{X=\{a\}\lor(X=\{b\}\lor X=\{a,b\})}{\FormulaRefAuto{TripleSetMembership}[thm]{12}}
  \proofstepwidestar[14]{X=\{a\}}{\rA}
  \proofstepwidestar[1,2,3,14]{X\in K}{\rIE{\FormulaRefAuto{a=b\vdash b=a}[thm]{14},7}}
  \proofstepwidestar[16]{X=\{b\}\lor X=\{a,b\}}{\rA}
  \proofstepwidestar[17]{X=\{b\}}{\rA}
  \proofstepwidestar[1,2,3,17]{X\in K}{\rIE{\FormulaRefAuto{a=b\vdash b=a}[thm]{17},8}}
  \proofstepwidestar[19]{X=\{a,b\}}{\rA}
  \proofstepwidestar[1,2,3,19]{X\in K}{\rIE{\FormulaRefAuto{a=b\vdash b=a}[thm]{19},9}}
  \proofstepwidestar[1,2,3,16]{X\in K}{\rOE{16,17,18,19,20}}
  \proofstepwidestar[1,2,3,10]{X\in K}{\rOE{13,14,15,16,21}}
  \nopagebreak
  \proofstepwidestar[1,2,3]{\PowNE{D}\subseteq K}{\FormulaRefAuto{SubsetIntroductionRule}[thm]{[10]\vdots 22}}
  \nopagebreak
  \proofstepwidestar[1,2,3]{K=\PowNE{D}}{\FormulaRefAuto{A\subseteq B,\, B\subseteq A\vdash A=B}[thm]{\rAEn{4},23}}
\end{tabproofwide}
